\documentclass[letterpaper,12pt]{article}
%\usepackage{harvard}
%\usepackage{subfigure}
%\usepackage[active]{srcltx}
\usepackage{graphicx}
\usepackage{pdflscape}
\usepackage{amsmath}
\usepackage{amsthm}
\usepackage{lscape}
\usepackage{enumerate}
\usepackage{float}
\usepackage{bm}
\usepackage{grffile}
\usepackage{relsize}
\usepackage{tablefootnote}
\usepackage{footnote}
\usepackage{sectsty}
\usepackage{xcolor}
\usepackage{caption}
\usepackage{color}
\usepackage{natbib}
\usepackage{multirow}
\usepackage{bbm}
\usepackage[T1]{fontenc}
\usepackage[multiple]{footmisc}
\usepackage{appendix}
\usepackage{times}
\usepackage{relsize}
\usepackage{caption}
\usepackage{newtxtext,newtxmath}
\usepackage{enumitem}


%%%%%%%%%%%%%%%%%%%%%%%%%%%%%%%%%%%%%%%%%%%%%%%%%%%%%%%%%%%%
%%%   Title Page Information
%%%%%%%%%%%%%%%%%%%%%%%%%%%%%%%%%%%%%%%%%%%%%%%%%%%%%%%%%%%%
\title{\vspace{-.5in} ECON712: Problem Set 1}
\author{Leandro Freylejer}
\date{Due date: February 13th, 2026 before 4pm }

%%%%%%%%%%%%%%%%%%%%%%%%%%%%%%%%%%%%%%%%%%%%%%%%%%%%%%%%%%%%
%%%   Double and Single Space Commands
%%%%%%%%%%%%%%%%%%%%%%%%%%%%%%%%%%%%%%%%%%%%%%%%%%%%%%%%%%%%
\newlength{\defbaselineskip}
\setlength{\defbaselineskip}{\baselineskip}
\newcommand{\setlinespacing}[1]{\setlength{\baselineskip}{#1 \defbaselineskip}}
\newcommand{\doublespacing}{\setlength{\baselineskip}{1.3 \defbaselineskip}}
\newcommand{\singlespacing}{\setlength{\baselineskip}{1\defbaselineskip}}
\renewcommand{\thesubsection}{\arabic{subsection}}
\newtheoremstyle{dotlessS}{}{}{\itshape}{}{\color{dg}\bfseries}{}{ }{}
\theoremstyle{dotlessS}
\newtheorem{theorem}{Theorem}
\newtheorem{lemma}{Lemma}
\newtheorem{proposition}{Proposition}
\newtheorem{assumption}[theorem]{Assumption}
\newtheorem{corollary}{Corollary}
\newtheorem{definition}{Definition}
\newenvironment{example}[1][Example]{\begin{trivlist}
\item[\hskip \labelsep {\bfseries #1}]}{\end{trivlist}}
\newenvironment{remark}[1][Remark]{\begin{trivlist}
\item[\hskip \labelsep {\bfseries #1}]}{\end{trivlist}}

\newcommand\Pitem{%
  \addtocounter{enumi}{-1}%
  \renewcommand\theenumi{\arabic{\enumi}'}%
  \item%
  \renewcommand\theenumi{\arabic{enumi}}%
}

%%%%%%%%%%%%%%%%%%%%%%%%%%%%%%%%%%%%%%%%%%%%%%%%%%%%%%%%%%%%
%%%   Set Margins Here
%%%%%%%%%%%%%%%%%%%%%%%%%%%%%%%%%%%%%%%%%%%%%%%%%%%%%%%%%
\textwidth=6.5in
\textheight=9in
\oddsidemargin=0.10in
\evensidemargin=0.10in
\topmargin=-0.5in
%\parskip=3pt

%%%%%%%%%%%%%%%%%%%%%%%%%%%%%%%%%%%%%%%%%%%%%%%%%%%%%%%%%%%%%
%%%   Actual Document Begins Here
%%%%%%%%%%%%%%%%%%%%%%%%%%%%%%%%%%%%%%%%%%%%%%%%%%%%%%%%%%%%%
\begin{document}
\pagenumbering{arabic}

\section*{Question IV: Empirical Application}

\begin{enumerate}[label=(\alph*)]
\item The following table replicates Table 2.2.1 in \textit{Mostly Harmless Econometrics} using the Labour Force Survey
\begin{table}[H]
\begin{center}
\caption{Comparison of individual characteristics BA vs. no-BA}
\label{table1}
\resizebox{.68\textwidth}{!}{%
\begin{tabular}{llll}
\hline
\hline
 Variable & No BA & BA & Diff \\
\hline
Ln(wage)     &       3.428***&       3.693***&       0.265***\\
                    &     (0.001)   &     (0.001)   &     (0.001)   \\
\hline
tenure (months)   &     100.670***&      91.100***&      -9.570***\\
                    &     (0.140)   &     (0.174)   &     (0.225)   \\
female          &       0.461***&       0.571***&       0.110***\\
                    &     (0.001)   &     (0.001)   &     (0.001)   \\
immigrant      &     0.181***&       0.334***&       0.152***\\
                    &     (0.001)   &     (0.001)   &     (0.001)   \\
Age 30-34    &       0.126***&       0.163***&       0.037***\\
                    &     (0.001)   &     (0.001)   &     (0.001)   \\
Age 35-39   &       0.130***&       0.167***&       0.037***\\
                    &     (0.001)   &     (0.001)   &     (0.001)   \\
Age 40-44   &       0.133***&       0.153***&       0.020***\\
                    &     (0.001)   &     (0.001)   &     (0.001)   \\
Age 45-49   &       0.130***&       0.130***&       0.001   \\
                    &     (0.001)   &     (0.001)   &     (0.001)   \\
Age 50-54   &       0.129***&       0.111***&      -0.018***\\
                    &     (0.001)   &     (0.001)   &     (0.001)   \\
Age 55-59   &       0.123***&       0.083***&      -0.041***\\
                    &     (0.001)   &     (0.001)   &     (0.001)   \\
Age 60-64   &       0.112***&       0.054***&      -0.058***\\
                    &     (0.001)   &     (0.000)   &     (0.001)   \\
\hline
\hline
\end{tabular}}
\end{center}
\captionsetup{font={scriptsize}}
\captionsetup{width=0.80\textwidth}
\captionsetup{skip=0pt}
\caption*{\textbf{Note:} This table displays average characteristics for individuals with and without a BA, defined as a four year University degree. The source are author's own calculations from the Canadian Labour Force Survey }
\end{table}
On average, workers with at least  four-year BA earn a statistically significant higher hourly wage. In particular, workers with four-year BA or higher earn 0.265 log points, or $30.34\%$ higher hourly wage. From the table it is clear that average workers across education levels are also different and that these differences are correlated with hourly wages. For instance, the proportion of females and immigrants are higher among workers with a BA, characteristics that are negatively correlated with earnings. On the other hand, tenure is, on average, lower for workers with a BA or higher, which correlates positively with earnings.

\item The following table replicates Table 2.2.2 in \textit{Mostly Harmless Econometrics} using the Labour Force Survey

\begin{table}[H]
\begin{center}
\caption{Returns to Finishing a four-year university degree}
\label{table1}
\resizebox{.82\textwidth}{!}{%
\begin{tabular}{llllll}
\hline
\hline
                    &\multicolumn{1}{c}{(1)}&\multicolumn{1}{c}{(2)}&\multicolumn{1}{c}{(3)}&\multicolumn{1}{c}{(4)}&\multicolumn{1}{c}{(5)}\\
\hline
4-year Degree     &       0.265***&       0.270***&       0.272***&       0.289***&       0.305***\\
                    &     (0.001)   &     (0.001)   &     (0.001)   &     (0.001)   &     (0.001)   \\
tenure         &               &               &       0.001***&       0.001***&       0.001***\\
                    &               &               &     (0.000)   &     (0.000)   &     (0.000)   \\
female        &               &               &               &      -0.151***&      -0.152***\\
                    &               &               &               &     (0.001)   &     (0.001)   \\
immigrant    &               &               &               &               &      -0.103***\\
                    &               &               &               &               &     (0.001)   \\
\hline
Age Group FE & No & Yes & Yes & Yes & Yes \\
Obs.                &    5.64e+05   &    5.64e+05   &    5.64e+05   &    5.64e+05   &    5.64e+05   \\
\hline
\hline
\end{tabular}}
\end{center}
\captionsetup{font={scriptsize}}
\captionsetup{width=0.80\textwidth}
\captionsetup{skip=0pt}
\caption*{\textbf{Note:} This table displays the returns to finishing a four-year degree or higher. The source are author's own calculations from the Canadian Labour Force Survey }
\end{table}
Adding controls increase the estimated effect of education on workers' earnings. From the previous table, it is clear that the estimated wage increases when controlling for female and immigrant as they are positively correlated with education; and they are negatively correlated with earnings. The effect of adding tenure in the regression is surprising as one will expect that, given the results of part (a), it will reduce the estimated coefficient. However, the table in part (a) does not includes age fixed effects when calculating differences in tenure across the two groups of workers (no BA vs. BA). Therefore, it could be the case that after controlling for age the tenure of workers with BA is larger. \\
\\
The coefficients estimates are not causal. There are many unobserved differences between workers with no BA and BA that are correlated with earnings, such as unobserved ability. Therefore, the previous regression suffers from omitted variable bias.  

\item Selection Bias with homogeneous effects of education is:
\begin{eqnarray*}
\text{Selection Bias}=\rho \cdot \left\{ \mathbb{E}[a_i|\bm{x}_i, Univ_i=1]-\mathbb{E}[a_i|\bm{x}_i, Univ_i=0] \right\}
\end{eqnarray*}

\end{enumerate}






\end{document}



